\section{Common fixed-field words and exact-width replay}

This appendix makes the local characteristic-two substitution explicit.
It does not import the complex convolution service from \cite{local}.

\paragraph{Newton symmetry over a ring.}
Let $P$ have as column $k$ the monomial coefficients of
$\prod_{a<k}(Z-u^a)$. It is unit upper triangular, and evaluation
gives $F_p(u)P=N$. With $D=\operatorname{diag}(N_{kk})$ and
$L=ND^{-1}$, this is the unit-lower/diagonal/unit-upper factorization
$F_p(u)=LDP^{-1}$. Such a factorization with invertible diagonal is
unique over a commutative ring: comparing two makes a unit lower
triangular matrix equal an upper triangular one, identifying the
diagonal and then the two unit triangular factors. Since $F_p(u)$ is
symmetric, its transposed factorization is another of that form.
Therefore $P^{-1}=L^T$ and $F_p(u)=ND^{-1}N^T$. The pivot inverses
are exactly the units proved in the main text; no field cancellation
has been assumed.

\paragraph{A shared eight-$H$ schedule.}
Write $D_a=\operatorname{diag}(a,1)$, with $a=\eta$ or $\eta^2$.
Direct multiplication in $\F_4$ gives
\[
 H^3=I,\qquad
 H D_\eta H^2 D_{\eta^2}H
   =\operatorname{diag}(\eta^2,\eta)J,
\]
and
\[
 W_4:=H D_\eta H D_{\eta^2}H D_\eta H D_{\eta^2}=I.
\]
Thus $I=W_4W_4$, $J$ is its displayed four-$H$ word with a leading
unit diagonal followed by $W_4$, and $H=H^4W_4$ all have eight $H$
letters with unit diagonals in common slots. Identity diagonals fill
missing gaps. The formulas include all identity pairs and never
select a zero scalar. Coefficients may depend on a pair label because
every intermediate letter preserves that label. Controlled
coefficients are evaluated only at the completed $K_T$ boundaries
described in the main proof.

\paragraph{Formal convolution over a coefficient algebra.}
For comparison, Cantor--Kaltofen \cite{cantor-kaltofen} give sharper
classical size bounds and logarithmic arithmetic depth. Their sparse
integer \emph{stage} matrices compose to a formal bilinear identity;
the composed matrices need not be sparse. We use the following
elementary triadic construction instead, so its descriptor generation
and coefficient types can be checked directly.

\begin{lemma}[Elementary formal convolution]
\label{lem:formal-convolution}
Over any commutative $\F_2$-algebra $A$, ordinary multiplication of
two length-$v$ polynomials has a division-free formal bilinear circuit
with size $O(v\log^8(2v))$, depth $O(\log(2v))$, binary fan-in,
and deterministic polynomial-$v$ wire-descriptor generation.
\end{lemma}
\begin{proof}
Use schoolbook products below the fixed threshold $v_0=6^6$, with
balanced output sums. Otherwise let $b$ be the least power of three
with $b^2\ge v$, put $s=\lceil v/b\rceil$, and let $t$ be the
least power of three with $t\ge2s-1$. Then
\[
 \sqrt v\le b<3\sqrt v,\quad s\le b,\quad
 t\le3b,\quad t<6s,\quad t\mid3b.
\]
Work formally in $R=A[Y]/(Y^{2b}+Y^b+1)$, free of rank $2b$
over $A$. The polynomial is $\Phi_{3b}$ reduced modulo two.
Its roots over $\overline{\F}_2$ have exact order $3b$: they
satisfy $Y^{3b}=1$, and none satisfies $Y^b=1$ since the defining
polynomial then equals one. Every proper divisor of $3b$ divides
$b$. Therefore $Y^j-1$ is coprime to this polynomial for
$0<j<3b$. The corresponding Bezout identity extends from $\F_2$
to $A$, proving principality of $Y$ and of
$\zeta=Y^{3b/t}$ at order $t$. The inverse DFT uses $\zeta^{-1}$;
its normalization is one because $t$ is odd.

Encode the input as $\sum_{i<s}A_i(Y)Z^i$, with
$\deg A_i<b$, and similarly for the other operand. Take their
$t$-point DFTs over $R$. At a frequency the value is \emph{one}
$R$ element, represented by $2b$ coefficients from $A$; no
unevaluated $Z$ variable remains. Multiply two such values by one
recursive ordinary length-$2b$ product over $A$, and reduce the
degree-$\le4b-2$ result modulo $Y^{2b}+Y^b+1$.
For $0\le r\le2b-2$, the high monomial reductions are
\[
 Y^{2b+r}=
 \begin{cases}
   Y^{b+r}+Y^r,&r<b,\\
   Y^{r-b},&b\le r\le2b-2.
 \end{cases}
\]
They cost $O(b)$ additions in constant depth.

The inverse DFT recovers
$C_j(Y)=\sum_{i+k=j}A_i(Y)B_k(Y)$ without wrap in $j$, since
$t\ge2s-1$. Each true $C_j$ has degree at most $2b-2$, below the
quotient degree, so reduction has lost no coefficient information.
Reconstruct the ordinary product as $\sum_j C_j(X)X^{bj}$:
adjacent length-$2b$ blocks overlap, and their overlaps are added,
not concatenated. Crop the known padded zero tail to length $2v-1$.

Multiplication by a root power in $R$ pads a $2b$-coordinate
vector to $3b$, rotates, and folds
$(a_0,a_1,a_2)\mapsto(a_0+a_2,a_1+a_2)$.
Each coordinate uses at most two inputs. A radix-three butterfly
therefore has $O(b)$ binary additions and constant depth over $A$.
Explicitly, at transform length $t=3u$ split inputs as $x_{3j+r}$,
$r=0,1,2$. Compute their three length-$u$ transforms $B_r(q)$
with root $\zeta^3$, then print
\[
 \widehat x_{q+\ell u}
   =\sum_{r=0}^2\zeta^{r(q+\ell u)}B_r(q),
 \qquad 0\le q<u,\quad \ell=0,1,2.
\]
This recursively specifies every butterfly and twiddle index.
The two forward and one inverse transforms, pointwise folds and
reconstruction give, for an effective absolute constant $C$,
\[
 S(v)\le tS(2b)+Ctb\log(2t),\qquad
 D(v)\le D(2b)+C\log(2t).
\]
These counts are in $A$ gates, not unit-cost $R$ operations.
Moreover
\[
 2bt<12bs\le12(v+b)\le24v,\qquad
 2b<6\sqrt v\le v^{2/3}\quad(v\ge v_0).
\]
Induction using $S(v)\le K v(\log_2v)^8$ above the direct branch
contracts the recursive term by
\[
 24(2/3)^8=6144/6561<1.
\]
The $O(v\log v)$ toll is absorbed by increasing a fixed $K$;
the finite direct branch fixes an effective base constant. This
proves the size bound. The logarithms of child lengths contract
by $2/3$, so the depth recurrence sums to $O(\log(2v))$.

All nonrecursive maps are the printed linear maps; recursive
leaves multiply a left-operand form by a right-operand form.
This proves a formal bilinear identity over $A$, without
$x^2=x$. For generation, compute $b,s,t$ by integer loops, emit
radix-three butterflies with rotation/fold source indices, and
recursively emit the $t$ children. Each descriptor has
$O(\log v)$ address bits; source indices can be obtained by
schoolbook arithmetic and sequential list scans in polynomial
time. This finite recursive printer is uniform, with no root
table, factorization, or coefficient-dependent branch.
\end{proof}

Here is a second explicit way to see that specialization leaves only
linear gates. After fixing the kernel, attach to each gate its
constant value $c_g$ and its homogeneous degree-one part $\ell_g$
in the remaining input indeterminates. An input $x_i$ has pair
$(0,x_i)$, a fixed scalar $a$ has $(a,0)$, and the rules are
\[
 (c,\ell)+(d,m)=(c+d,\ell+m),\qquad
 (c,\ell)(d,m)=(cd,c\,m+d\,\ell).
\]
The constants are prepared once. Only additions and multiplication
by these fixed constants remain in the linear DAG. The rules are
formal truncation modulo total degree at least two; they use no
Boolean relation $x_i^2=x_i$. Since fixed-kernel convolution is
homogeneous linear, its output is exactly the degree-one part.
The transformation changes size and depth by constant factors.

Bounded-degree copy trees, including copies of inputs and outputs,
cost size proportional to the number of use edges. Replacing each
gate's fan-out by its balanced tree adds at most $O(\log v)$ depth
per original circuit layer. The deliberately loose linear-DAG bounds
are therefore
\[
 g(v)=O(v\log^8(2v)),\qquad H(v)=O(\log^2(2v)).
\]
All wire and scalar descriptors have polynomial-time generation in
$v$ and the explicit coefficient dimension. Kernel-specific constant
propagation is charged once in that setup, with the actual coefficient
words retained. This construction is
independent of the faster multiplier being proposed.

\paragraph{Replay on arbitrary dirty coordinates.}
Let a linear DAG implement $M:A^e\to A^a$, and assign a distinct
borrowed coordinate to each of its $g$ gates. The physical source
coordinates $x$, target coordinates $y$, and borrowed coordinates $z$
are pairwise disjoint. Every output is a designated wire; any additions
needed to form it and any input/output copy trees are included in the
counted DAG. Thus the output-read edges also have bounded degree.
In the notation below, $F,D$ name these actual port reads, rather than
uncosted dense matrix applications.
Replace gate assignment
by addition of its predecessor form into that coordinate, in
topological order. For initial borrowed vector $z$ this sweep has
form $z\mapsto Uz+Tx$, where $U$ is unit lower triangular.
If the prescribed reads are $Fz+Dx$, zero initialization would compute
the DAG, so $FT+D=M$.
Sweep, add these reads to $y$, and undo the sweep. The net increment
is $Mx+FUz$, while $z$ is restored. Now sweep with every direct edge
from $x$ omitted, including edges to input-copy-tree roots, and add
only the borrowed-coordinate read $FUz$ to $y$,
and undo again. In characteristic two the two $FUz$ terms cancel.
This realizes $(x,y,z)\mapsto(x,y+Mx,z)$ for every initial $z$,
including coordinates correlated with current sources and targets.

At each DAG depth destinations are fresh gate coordinates and
predecessors are earlier ones. Bounded fan-in and fan-out give a
bounded-degree graph of shear edges. A constant edge coloring makes
each color disjoint pairs. Apply its colors before advancing to the
next depth; undo in reverse depth and color order. The output-port
edges into the distinct $y$ coordinates have the same constant-color
bound. In the second read omit every direct $x$-to-output edge as well.
Zero-coefficient edges may be omitted throughout.
The four sweeps and two output rounds
thus use $O(H+1)$ shear layers, on exactly $e+a+g$ coordinates.
No clean gate values or fresh zero coordinates are assumed.
Other parent coordinates can be selected as borrowed gates, even if
they contain data from an earlier chunk operation; the replay restores
them exactly. A source/target alias is outside this contract.

\paragraph{Explicit rectangular displacement reconstruction.}
\begin{lemma}[Exact-width unit Toeplitz word]
\label{lem:exact-width-toeplitz}
Over any commutative $\F_2$-algebra $A$, a unit lower triangular
Toeplitz matrix $T_v$ has a word on exactly its original $v$ coordinates
of $O(\log^{19}(2v))$ layers. Each layer consists of shears
$(x,y)\mapsto(x,y+\mu x)$ on disjoint ordered pairs.
No clean coordinate is required. Descriptors and coefficients can be
prepared by polynomially many operations on indices, finite lists,
and public coefficient-algebra elements.
\end{lemma}
\begin{proof}
Write $(T_v)_{ij}=h_{i-j}$ for $i\ge j$, with $h_0=1$.
Its reciprocal sequence is prepared by $g_0=1$ and
$g_k=-\sum_{j=1}^k h_jg_{k-j}$; no coefficient inversion is needed.
For $v\ge2$ split at $s=\lfloor v/2\rfloor$ and put $r=v-s$.
The matrix has blocks
\[
 T_v=\begin{pmatrix}T_s&0\\B&T_r\end{pmatrix}.
\]
After applying the two diagonal children in place, a cross update
$E=BT_s^{-1}$ gives this matrix exactly. Its entries are
\[
 E_{ij}=\sum_{\nu=j}^{s-1}h_{i-\nu}g_{\nu-j}.
\]
For consecutive row/column chunks, the interior difference is
\[
 E_{ij}-E_{i-1,j-1}=-h_{i-s}g_{s-j}.
\]
Let $M$ be that rectangular chunk and let $J_a,J_e$ be lower shifts.
Its difference $\mathcal D=M-J_aMJ_e^T$ is the displayed outer
product plus first-row and first-column corrections, hence the sum
of three explicit outer products. More explicitly, let the chunk's
absolute row and column indices start at $i_0\ge s$ and $j_0<s$,
and use zero-based local indices. Set
\[
 v_i=-h_{i_0+i-s},\qquad w_j=g_{s-j_0-j},\qquad
 r_j=M_{0j}-v_0w_j,
\]
and set $c_0=0$, $c_i=M_{i0}-v_iw_0$ for $i>0$. Then, with $e_0$
the first standard vector of the appropriate length,
\[
 \mathcal D=vw^T+e_0r^T+ce_0^T.
\]
The interior follows from the displayed difference identity, the
second term corrects the complete first row, and the third corrects
the remaining first column without counting $(0,0)$ twice. Every
entry is given by finite sums and products of the prepared Toeplitz
and reciprocal coefficients. This uses no field rank test.
Nilpotence gives
$M=\sum_{l\ge0}J_a^l\mathcal D(J_e^T)^l$.
For one term $vw^T$, its action on $x$ is
\[
 b_l=\sum_{j=l}^{e-1}w_{j-l}x_j,\qquad
 c_i=\sum_{l=0}^{e-1}v_{i-l}b_l.
\]
Here $v_k=0$ outside $0\le k<a$.
For literal convolution windows put
$W(Z)=\sum_{j<e}w_jZ^j$, $X^{\rm rev}(Z)=\sum_{j<e}x_jZ^{e-1-j}$
and $V(Z)=\sum_{i<a}v_iZ^i$. The polynomial with coefficients $b_l$ is
\[
 \operatorname{rev}_e\bigl((W X^{\rm rev})\bmod Z^e\bigr),
\]
and the output polynomial is its product with $V$, reduced modulo $Z^a$.
Here $\operatorname{rev}_e P=Z^{e-1}P(Z^{-1})$ at fixed padded width.
The two ordinary products have output lengths at most $2e-1$ and
$a+e-1$, respectively. If $a,e\le b$, both are at most $2b-1$.
Their padding zeros are formal constants; they introduce no extra
runtime source coordinates. Reversal and selection specify DAG ports.
The three displacement terms therefore use six fixed-kernel convolution
circuits and their sum, with total gate bound $O(b\log^8(2b))$ and
depth $O(\log^2(2b))$ including all intermediate and copy coordinates.

At setup enumerate $b=1,\ldots,v$, partition both sides of the cross
block into consecutive chunks of width at most $b$, and print the
actual DAG for every pair, including ragged final chunks and its
input/output copies. Choose the largest $b$ for which
$e+a+g\le v$ for every pair. For a uniform fixed $C$ the count is
$e+a+g\le2b+Cb\log^8(2b)$. Thus
$b_0=\lfloor c v/\log^8(2v)\rfloor$ passes for sufficiently large $v$
and a sufficiently small fixed $c>0$; the chosen size is at least $b_0$.
If none passes, use the direct triangular shear word, updating targets
in descending order from their earlier source coordinates.
This branch occurs only in a fixed finite range.

For each accepted pair its source chunk lies in $[0,s)$ and its target
chunk in $[s,v)$, hence they are disjoint. Choose the $g$ borrowed
coordinates from their complement inside this same $v$-coordinate
parent. Replay restores that complement before the next pair.
The source half is unchanged by every completed cross update, so the
sum of all chunk updates is exactly $E$.
There are
$O(\log^{16}v)$ chunk pairs, each with $O(\log^2v)$ replay layers.
The diagonal children use only their own disjoint intervals.
Run their matching layers in parallel, padding the shorter word
with identity layers. Their union is again a matching. Thus
\[
 d(v)\le\max(d(\lfloor v/2\rfloor),d(\lceil v/2\rceil))
       +O(\log^{18}(2v)),\qquad d(v)=O(\log^{19}(2v)).
\]
All coordinates belong to the original $v$-coordinate parent. Preparing
and counting graphs, reciprocal coefficients, outer-product generators
and edge colors uses only polynomial-size finite lists and polynomially
many public coefficient operations. The search is setup for the public
kernel, never a per-payload-record computation.
\end{proof}

The Newton Toeplitz sequence is $h_j=H_j^{-1}$, so $h_0=1$ as required
by Lemma~\ref{lem:exact-width-toeplitz}. Transposing a layer word reverses
its layer order and each shear's source/target direction; this supplies
the upper triangular word on the same coordinates at the same layer
count. This transposition is on $A$-coordinate indices and retains each
coefficient $\mu$. It is a descriptor operation, not a physical array transpose.
Consequently $F_p=ND^{-1}N^T$ uses two such Toeplitz words and a fixed
number of unit-diagonal stages. Synchronize different prime axes by
identity matching layers. With $p\le R_0$, preparing all their words and
actual $B$-coefficient tables has bit cost polynomial in $R_0,D',L_0$,
as in the coefficient setup below. Each matching layer is then replaced
by the fixed local slot word in the main text; these abstract layer
counts do not assign unit bit cost to growing scalars or tape movement.

\paragraph{The actual constant extension.}
Every irreducible factor of $\Phi_{\alpha'}$ has degree
$r=\operatorname{ord}_{\alpha'}(2)$, so $S$ is a product of copies
of $\F_{2^r}$. A factor tensored with $\F_4$ splits into
$\gcd(r,2)$ fields of degree $r_B=\operatorname{lcm}(r,2)$.
Therefore the same Frobenius exponent in the main proof selects
$\mu=1$ in every field component of $B$, with $r_B\le2D'$.
Computing $r$ by repeated multiplication by $2$ modulo $\alpha'$
needs at most $D'$ iterations and has polynomial small setup cost.
Exponentiation for the selector needs $O(r_B)$ ring products.
The fixed $\F_{16}$ codec inside the binary $H$ supplier need not
embed into $B$: the completed $H$ acts identically on each of the
$2D'$ binary lanes and is thus the required $B$-symbol map.

The central $\F_4$ in $B$ and a fixed field codec used inside a
binary routing/H kernel need not have the same intermediate bit
encoding. The replacement is made only at a completed binary kernel:
its target matrix is $H$, $I$ or $J$ on whole symbols and acts
identically on every binary lane. It therefore implements the required
$B$-linear map in the caller's original representation. Outer
multiplications by elements of $B$ are inserted only after that
boundary and always use the separately charged coefficient-ring
arithmetic. Even a central $\F_4$ action on two concatenated $S$
lanes can be implemented by a constant number of aligned word scans;
no growing basis conversion or free bit transpose is assumed.

\paragraph{Actual growing-word arithmetic and unit setup.}
\begin{lemma}[Fixed monic remainder at any input width]
\label{lem:fixed-monic-remainder}
Fix a monic binary polynomial $f$ of degree $d\ge1$ and a public input
width $r\ge d$. After polynomial-$r$ setup, the remainder modulo $f$
of any binary polynomial of degree $<r$ costs $O(M_{\rm old}(r))$
bit steps on the fixed tape model, including all scans and returns.
\end{lemma}
\begin{proof}
If $r=d$, the input is already reduced. Otherwise set $k=r-d>0$.
The reversal $f^*=X^df(X^{-1})$ has constant coefficient one by
monicity. Prepare $h=(f^*)^{-1}\bmod X^k$ once, for example by its
coefficient recurrence with elementary polynomial-time tape arithmetic.
For the dividend $w$ put
\[
 w^*=X^{r-1}w(X^{-1}),\qquad
 q^*=(w^*h)\bmod X^k,\qquad
 q=X^{k-1}q^*(X^{-1}),\qquad R=w-qf.
\]
All reversals have the displayed fixed padded widths. To prove the
formula, write the unique monic division $w=Qf+R_0$ with
$\deg Q<k$ and $\deg R_0<d$. Reversal gives
$w^*=Q^*f^*+X^{r-1}R_0(X^{-1})$; the second term is divisible by
$X^k$. Multiplying by $h$ modulo $X^k$ therefore recovers $Q^*$,
so $q=Q$ and $\deg R<d$.

The first product may truncate $w^*$ to its first $k$ coefficients.
It has two inputs of length $k$; the other product has input lengths
$k$ and $d+1$, each at most $r$. Thus two calls to the independent
word supplier cost $O(M_{\rm old}(r))$, and all reversals, cropping,
XOR and head returns add $O(r)$. There is no per-dividend Euclidean
algorithm. The prepared inverse length is $k$, even when $d\ll r$.
\end{proof}

Put $d=D'$ and $f=\Phi_{\alpha'}$. Store a $B$ symbol as the two
consecutive $d$-bit $S$ words in $a+b\eta$. If
$P=ac$, $Q=be$, and $T=(a+b)(c+e)$ in $S$, then
\[
 (a+b\eta)(c+e\eta)=(P+Q)+(T+P)\eta.
\]
Three binary polynomial products modulo fixed $f$, supplied by
Lemma~\ref{lem:old-word}, with linear XOR scans, therefore suffice.
Their degree-$\le2d-2$ outputs use
Lemma~\ref{lem:fixed-monic-remainder} at input width $2d-1$.
Its reciprocal is prepared at length $d-1$ for $d\ge2$; $d=1$
is direct. This proves $O(M_{\rm old}(d))$ application time with
no dense $d$-by-$d$ matrix operation.
The runtime word supplier is classical integer SS with carry-free
binary packing; it is separate from the formal characteristic-two
convolution circuit used to construct the local linear DAG above.

The involution $\eta\mapsto\eta+1$ has norm
$N(a+b\eta)=a^2+ab+b^2\in S$. For a unit $x=a+b\eta$, its norm is
an $S$ unit: it is a product of units in $B$, and an inverse in
the free extension $B=S\oplus S\eta$ supplies an inverse in $S$
by taking its constant coefficient. Thus
\[
 x^{-1}=((a+b)+b\eta)\,N(x)^{-1}.
\]
Polynomial extended Euclid with $f$ suffices for the norm inverse,
without factoring $B$. Such inverses are prepared once in polynomial
time in the small dimensions. The selector exponent $2^{r_B}-1$
also needs only $O(r_B)$ binary squaring/product steps, with $r_B\le2d$.
For example,
from $z=\mu+1$ repeat $z\gets z^2(\mu+1)$ exactly $r_B-1$ times.
No possibly nonunit $\mu+1$ is inverted. The one-time polynomial
setup is separate from every runtime scalar application.

\paragraph{Scalar lists on fixed tapes.}
At a synchronized local slot, each axis has a public table
$a_i(x_i)\in B^*$ in that slot's current coordinate order; the
tensor diagonal is $\prod_i a_i(x_i)$. Pairing and coefficient
choices are axis-local. Jointly controlled CRT coefficients use
their different generator and are not covered by this factorization.

Let $M_j=\prod_{i\le j}n_i$, $M_0=1$, with axes in physical
lexicographic order. Start with the one-word list $[1]$.
For every partial product $c$ on the current list, scan the entire
next local table and print
$ca_j(0),\ldots,ca_j(n_j-1)$. Rewind that local table, then advance
to the next $c$. Its rewind costs the size of the child block just
printed. Alternate two expansion tapes and erase the previous list
at the end of each level. This keeps the produced list in the
current native order and uses a fixed tape set.

Level $j$ uses $M_j$ growing-ring products and
$O(M_j(d+L))$ copying, counters and returns. Here $2d\ge L$, and
all $n_i\ge2$, so
\[
 \sum_j M_j<2M,\qquad
 \cost(\hbox{generate and apply one tensor diagonal})
       =O(MM_{\rm old}(d)).
\]
Apply the final list to the payload stream in lockstep and erase it.
Prepare or select each local table once per axis/slot in the
polynomial small setup, using one shared archive and one reusable
active-table tape; there is no tape per axis or indexed archive lookup per
payload record. Rebuild the list from the next slot's tables after
completed relabeling, $H$ and $J$ calls. Transporting an old scalar
list and presuming alignment is not part of the program.
